\begin{explanationbox}
	\textbf{\textcolor{CExplanation}{一句话直觉}}

	四点共圆等价于四边形对角互补，这是圆周角内接四边形的直观体现。

	\medskip
	\textbf{\textcolor{CExplanation}{核心拆解}}
	\begin{description}[style=nextline, leftmargin=2.8em, labelsep=0.5em, itemsep=0.45em]
		\item[\textbf{要点一（充要条件）：}] 对角互补是四点共圆的充要条件，具有双向性。
		\item[\textbf{要点二（外角转化）：}] 利用外角等于内对角，将圆内角关系转化为直线角，简化图形。
		\item[\textbf{要点三（同侧等角）：}] 对同底 $AB$ 的 $\angle ADB=\angle ACB$ 可避开四边形，直接利用三角形。
	\end{description}

	\medskip
	\textbf{\textcolor{CExplanation}{几何本质（圆周角定理）}}

	圆内接四边形的一组对角所对的两段弧恰好合成整个圆，由于圆周角等于所对弧度数的一半，故两角之和为 $180^\circ$。

	\medskip
	\textbf{\textcolor{CExplanation}{代数意义（角度方程）}}

	设定相关角，通过三角形内角和与列方程，可将几何条件转化为角度等式 $\alpha+\gamma=180^\circ$，便于代数运算。

	\medskip
	\textbf{\textcolor{CExplanation}{考点价值（四点共圆判定）}}
	\begin{itemize}[leftmargin=*, itemsep=0.5em, topsep=0.4em]
		\item \textbf{考法一（对角求和）：} 直接计算四边形两个对角，若和为 $180^\circ$ 则判定共圆。
		\item \textbf{考法二（等角转换）：} 利用外角或同侧角相等间接判定共圆。
	\end{itemize}

	\medskip
	\textbf{\textcolor{CExplanation}{顿悟点}}

	将四边形补全为圆，对角互补与圆周角对应关系是转化的桥梁。

	\medskip
	\textbf{\textcolor{CExplanation}{使用场景}}
	\begin{itemize}[leftmargin=*, itemsep=0.5em, topsep=0.4em]
		\item \textbf{场景一（角度证明）：} 已知夹角度数或比例，判断四点共圆后引出等角关系。
		\item \textbf{场景二（图形简化）：} 在外角或同侧角相等情形下，直接判定共圆而无需完整四边形。
	\end{itemize}
\end{explanationbox}
